
\chapter{Part III Review: Invariance, Amplitudes, and Oscillators}

Relativity identifies the invariant spacetime interval.  Quantum mechanics identifies states as complex amplitude vectors.  The harmonic oscillator supplies the algebraic unit that will later appear once for every field mode.

\section{A guided reconstruction}

% QFTFIGURE-BEGIN partreview03-01-part-map
\qftfigure{figures/instructional/partreview03/partreview03-01-part-map.pdf}{The map collects the part's main constructions into a visual route so the reader can see how the chapters support one another.}{fig:partreview03-01-part-map}
% QFTFIGURE-END partreview03-01-part-map












The mass shell \(p_\mu p^\mu=m^2\) and the commutator \([q,p]=\ii\) are both compact statements with long consequences.  The first constrains the form of relativistic wave equations; the second produces ladder operators and a discrete spectrum:
\[
  H=\omega\left(a^\dagger a+\frac12\right).
\]
Together they prepare the statement that a free relativistic field is a collection of quantum oscillators labelled by momentum.

\begin{exercise}
Explain why the oscillator spectrum is needed before the particle interpretation of a free field can be built.
\begin{answer}
Each momentum mode of the free field is an oscillator, and its occupation number becomes the number of particles in that mode.
\end{answer}
\end{exercise}
